% !TEX root = ../main.tex
\chapter{\cmm{}: fundamentos}
\label{cap:cmm-fundamentos}

A \cmm{} (lê-se \enquote{C mais-menos}; nos arquivos, extensão \file{.cmm}) é a linguagem principal do SAPHO: um dialeto de C criado pelo NIPS-CERN, enxuto o bastante para virar hardware eficiente e expressivo o bastante para algoritmos de processamento de sinais --- incluindo números complexos como tipo nativo e álgebra linear em notação de Dirac. Este capítulo apresenta a estrutura de um programa; os capítulos seguintes cobrem tipos, operadores, E/S, biblioteca e recursos avançados.

\begin{nota}
Tudo o que está nesta Parte foi extraído da gramática oficial do compilador (\file{CMMComp.y} e \file{CMMComp.l} do repositório \repo{yanc}, toolchain 5.3). Quando este manual diz que algo \enquote{não existe} na linguagem, é porque não existe na gramática --- o compilador rejeitará o código.
\end{nota}

\section{A estrutura de um programa}

Um programa \cmm{} é uma sequência de três tipos de elemento, em qualquer ordem: \textbf{diretivas} (linhas iniciadas por \code{\#}), \textbf{declarações de variáveis globais} e \textbf{funções}. O ponto de entrada é obrigatório: uma função \code{void main()}. Sem ela, o compilador encerra com erro.

A forma canônica de um programa SAPHO é: diretivas no topo (configurando o processador), depois o \code{main()} com um laço infinito \code{while (1)} --- afinal, o programa modela um circuito que roda continuamente, processando amostras para sempre:

\begin{lstlisting}[style=cmm]
#PRNAME media_movel
#NUBITS 16
#NDSTAC 4
#SDEPTH 2
#NUIOIN 1
#NUIOOU 1
#NBMANT 10
#NBEXPO 5

void main()
{
    int x[4];    // historico das ultimas 4 amostras
    int soma;

    while (1)
    {
        x[3] = x[2];          // desloca o historico
        x[2] = x[1];
        x[1] = x[0];
        x[0] = in(0);         // le nova amostra da porta 0

        soma = x[0] + x[1] + x[2] + x[3];
        out(0, soma >> 2);    // media = soma/4, sem usar divisor
    }
}
\end{lstlisting}

Este é o \code{media_movel}, o processador que acompanha o manual inteiro: um filtro de média móvel de 4 amostras. Repare em duas escolhas típicas de quem escreve para o SAPHO: o histórico é um array deslocado à mão, e a divisão por 4 é um deslocamento de bits (\code{>> 2}) --- assim o processador gerado não precisa instanciar um divisor, que é um dos blocos mais caros da ULA.

\section{Diretivas: configurando o processador no próprio código}
\label{sec:diretivas}

As diretivas são linhas que começam com \code{\#}, uma por linha, \textbf{sem ponto e vírgula}. Elas definem o nome e a arquitetura do processador que será gerado. Quando você cria um processador pela AURORA (Capítulo~\ref{cap:criando-processadores}), o formulário preenche essas diretivas por você no \file{.cmm} inicial; depois disso, o arquivo é a fonte da verdade --- editar a diretiva muda o processador na próxima compilação.

\begin{longtable}{@{}llp{0.52\linewidth}@{}}
\toprule
\textbf{Diretiva} & \textbf{Padrão} & \textbf{Significado} \\
\midrule
\endhead
\code{\#PRNAME <nome>} & --- & Nome do processador (deve casar com o nome usado no projeto). \\
\code{\#NUBITS <n>} & 23 & Largura da palavra em bits (inteiros em complemento de dois). \\
\code{\#NBMANT <n>} & 16 & Bits de mantissa do ponto flutuante. \\
\code{\#NBEXPO <n>} & 6 & Bits de expoente do ponto flutuante. \\
\code{\#NDSTAC <n>} & 10 & Profundidade da pilha de dados. \\
\code{\#SDEPTH <n>} & 10 & Profundidade da pilha de sub-rotinas (limita chamadas aninhadas). \\
\code{\#NUIOIN <n>} & 1 & Número de portas de entrada. \\
\code{\#NUIOOU <n>} & 1 & Número de portas de saída. \\
\code{\#NUGAIN <n>} & 64 & Divisor fixo usado pela função \code{norm()}. \\
\code{\#FFTSIZ <n>} & 8 & Tamanho da FFT ($2^n$ pontos) para o endereçamento bit-reverso. \\
\bottomrule
\end{longtable}

\begin{aviso}
Lembre da restrição estrutural: \texttt{NUBITS} = \texttt{NBMANT} + \texttt{NBEXPO} + 1. No exemplo acima, $16 = 10 + 5 + 1$. A AURORA valida isso ao criar o processador, mas se você editar as diretivas à mão, é sua responsabilidade manter a igualdade.
\end{aviso}

Existem ainda duas diretivas \emph{comportamentais}, que aparecem no meio do código (não no topo) e são explicadas no Capítulo~\ref{cap:cmm-avancado}: \code{\#PRACA} (ponto de retorno da interrupção) e \code{\#TOAQUI} (marcador refletido no pino \code{cheguei}).

\section{Constantes com \texttt{\#define}}

A \cmm{} tem um pré-processador mínimo embutido: \code{\#define NOME corpo} define uma constante simbólica, expandida onde o nome aparecer.

\begin{lstlisting}[style=cmm]
#define SCALE  3
#define OFFSET 7

void main()
{
    int x;
    while (1)
    {
        x = in(0);
        x = x + SCALE;
        out(0, x);
        x = x + OFFSET;
        out(0, x);
    }
}
\end{lstlisting}

Três limitações importantes em relação ao C tradicional:
\begin{itemize}
  \item só existem \emph{defines de objeto} --- macros com argumentos, como \code{\#define DOBRO(x) ((x)*2)}, não são suportadas;
  \item não há \code{\#include} nem \code{\#ifdef} no lado \cmm{} (o fluxo C++, sim, tem pré-processador completo --- Capítulo~\ref{cap:cmm-avancado});
  \item o limite é de 256 defines por programa.
\end{itemize}

O editor da AURORA reconhece os seus \code{\#define} e os colore como constantes conforme você digita.

\section{Comentários}

Como em C: \code{//} comenta até o fim da linha, \code{/* ... */} comenta blocos.

\section{Onde o arquivo vive}

Dentro de um projeto AURORA, o código de cada processador fica em \file{<projeto>/<processador>/Software/<processador>.cmm}. A compilação gera o assembly em \file{Software/}, o Verilog e as imagens de memória em \file{Hardware/} e o \emph{testbench} em \file{Simulation/} (Capítulo~\ref{cap:compilacao}).
